%% P21 --- Optymalizacja stochastyczna i różniczkowanie losowania
%%
%% Bliźniak programs/en/P21-stochastic-optimisation.tex. Ramka w ramkę, makro
%% w makro, liczba w liczbę; tools/parity.py porównuje oba pliki.
%%
%% UWAGA DLA TŁUMACZA: C4, C8 i C12 porównują znaczniki strukturalne, wzory i
%% literały liczbowe PO KOLEI. Cyfra zostaje cyfrą, odwołanie do programu
%% prowadzi zdanie zamiast stać za wzorem, a skrót jest częścią strumienia
%% znaczników. \enquote{batch} tłumaczy się jako \enquote{partia}, ale
%% \enquote{estymator}, \enquote{wariancja} i \enquote{gradient} zostają.

\program[Optymalizacja stochastyczna]{Optymalizacja stochastyczna i
różniczkowanie losowania}
\label{prog:P21}
\index{optymalizator!stochastyczny}
\index{gradient!z partii}

\begin{outcomes}
\outcome{1--7}{powiedzieć, co szacuje gradient z partii i co
  \emph{nieobciążoność} obiecuje, a czego nie}
\outcome{8--12}{powiedzieć, jak szum zależy od rozmiaru partii i dlaczego jest
  to zależność pierwiastkowa}
\outcome{13--17}{powiedzieć, dlaczego w ogóle zdarza się ogromny krok i czego
  dotyczy decyzja o progu obcinania}
\outcome{18--22}{podać, co utrzymuje stałym reguła liniowa, a co pierwiastkowa,
  i powiedzieć, która jest ludową mądrością}
\outcome{23--25}{powiedzieć, co ukrywa wygładzona krzywa straty i o ile}
\outcome{26--30}{powiedzieć, kiedy akumulacja mikropartii nie jest jedną dużą
  partią i dlaczego dzieje się to po cichu}
\outcome{31--39}{nazwać dwa sposoby na gradient przez krok losowania i
  powiedzieć, co rozstrzyga między nimi}
\end{outcomes}

\begin{quiz}
\item Gradient z partii nazywa się \emph{nieobciążonym}. Co to obiecuje?
  \teachesat{3--4}
  \answerto{Że uśrednienie po każdej partii, jaka mogła zostać wylosowana, daje
  dokładnie pełny gradient. Nie obiecuje zupełnie nic o partii, którą
  wylosowano.}
\item Zwiększasz rozmiar partii szesnastokrotnie. O jaki czynnik spada szum w
  gradiencie? \teachesat{10--11}
  \answerto{O cztery. Wariancja spada jak $1/B$, więc rozrzut \dash{} czyli to,
  co widać \dash{} spada jak $1/\sqrt{B}$. Ten pierwiastek jest źródłem
  malejących zysków i jest dokładny, a nie przybliżony.}
\item Zmniejszasz partię o połowę. Co powinno stać się z krokiem?
  \teachesat{18--19}
  \answerto{Zależy, co chcesz utrzymać stałe, a dwie popularne reguły
  utrzymują różne rzeczy. Zmniejszenie kroku razem z partią zachowuje
  \emph{stosunek}, w którym reguły są zwykle podawane; podzielenie go przez
  $\sqrt{2}$ zachowuje szum w aktualizacji. Żadna z nich nie jest prawem.}
\item Akumulowane są cztery mikropartie. Czy to to samo, co jedna partia
  czterokrotnie większa? \teachesat{26--27}
  \answerto{Tylko wtedy, gdy każda mikropartia zawiera tyle samo prawdziwych
  tokenów. W przeciwnym razie akumulowana strata jest średnią ze średnich, a
  strata dużej partii jest średnią z puli, i są to dwa różne cele.}
\item Dwa estymatory gradientu przez krok losowania są oba nieobciążone. Co
  jeszcze pozostaje do wyboru między nimi? \teachesat{33--34}
  \answerto{Ich wariancja, a może różnić się o rzędy wielkości. Bycie
  poprawnym średnio to nie to samo, co bycie użytecznym, i przy stu wymiarach
  jeden z dwóch jest wciąż widocznie błędny po czterdziestu tysiącach próbek.}
\end{quiz}

% ============================================================
\section{Gradient, o który nie prosiłeś}
\label{sec:P21-estimator}
\index{gradient!jako estymator}

\begin{fr}
Program~\ref{prog:P20} rozłożył optymalizator na części i ani razu nie zapytał,
skąd bierze się gradient. Każdy marsz w nim używał gradientu całej straty, a
żaden przebieg treningowy nigdy tego nie zrobił.

To, co przebieg liczy, to gradient straty na \emph{partii}. Powiedz, czym to
czyni liczbę, wzdłuż której kroczysz.
\dotline
\nextframe
\end{fr}

\begin{fr}
\begin{ansblock}
\textbf{Oszacowaniem} gradientu, którego chciałeś, zbudowanym z przykładów,
które akurat zostały wylosowane.
\end{ansblock}

To cały temat tego programu i warto powiedzieć to wprost, zanim cokolwiek
zostanie wyprowadzone: \textbf{szum nie jest wadą metody, szum jest metodą}.
Nikogo nie stać na pełny gradient i nikt go nie chce \dash{} trening robi więc
zaszumiony krok bardzo wiele razy, a własności tego szumu rozstrzygają niemal
wszystko, na co ludzie narzekają.

\mermaidfig{p21-one-draw}{Jedno losowanie i zbiór, z którego pochodzi.}{jedno
losowanie, jeden zbiór}
\end{fr}

\begin{fr}
Zacznij od tego, co jest obiecane. Weź populację $\val{p21.pop.n}$ liczb
oznaczających gradienty na przykład i każdą możliwą partię
$\val{p21.pop.b}$ z nich \dash{} jest ich $\val{p21.pop.subsets}$.

Uśrednij $\val{p21.pop.subsets}$ średnich z partii. Co otrzymasz?
\nextframe
\end{fr}

\begin{fr}
\ans{$\val{p21.pop.mean}$}
Własną średnią populacji, i to \textbf{dokładnie}: skrypt robi całość na
ułamkach, więc nie ma w tym żadnej tolerancji. Tym właśnie jest
\emph{nieobciążoność} i na tym opiera się każda metoda z partiami.

\begin{note}
Każda partia została wylosowana, więc jest to \textbf{dowód} dla tej populacji,
a nie świadectwo o niej, co jest rozróżnieniem Programu~\ref{prog:P14}. Na
prawdziwym zbiorze treningowym to samo zdanie jest twierdzeniem o losowaniu, a
nie eksperymentem, a arytmetyka jest identyczna.
\end{note}

Teraz połowa, która znaczy więcej. Te $\val{p21.pop.subsets}$ średnich dają w
średniej dokładnie $\val{p21.pop.mean}$. Zgadnij, jak daleko od niej jest
najgorsza z nich.
\dotline
\nextframe
\end{fr}

\begin{fr}
\begin{ansblock}
Biegną od $\val{p21.pop.lo}$ do $\val{p21.pop.hi}$.
\end{ansblock}

\transcript{p21-every-batch}

Rozpiętość $\val{p21.pop.span}$ wokół wielkości $\val{p21.pop.mean}$, a jedna z
nich ma zły znak.
\end{fr}

\begin{fr}
\begin{trapbox}
\textbf{\enquote{Gradient jest nieobciążony, więc mogę mu ufać.}}
Nieobciążoność to stwierdzenie o \emph{zbiorze}: uśrednij po każdej partii,
jaka mogła zostać wylosowana, a wynik jest dokładny. Nie mówi zupełnie nic o
partii, którą masz przed sobą, a tutaj partie otaczają odpowiedź z obu stron i
jedna z nich wskazuje w przeciwną stronę.

Dwa słowa, które chodzą w parze, to \emph{nieobciążony} i \emph{wariancja}, a
metoda mająca pierwsze bez kontroli nad drugim to metoda, której pojedyncze
kroki są szumem.
\end{trapbox}
\end{fr}

\begin{fr}
\begin{rigourbox}
\textbf{Ten program używa wariancji, zanim książka ją zdefiniowała}, i jest to
celowe, a nie przeoczenie. Czym jest zmienna losowa, czym wariancja i jak
średnie się koncentrują, to Programy \ref{prog:P24} i \ref{prog:P25}, dwie
części później; notatki programowe zapisują, dlaczego kolejność jest taka i
dlaczego przestawienie części zepsułoby sześć innych zależności.

Potrzebujesz tu jednego zdania i jest to zdanie, które i tak byś zgadł:
\textbf{wariancja} wielkości to średnia z kwadratu jej odległości od własnej
średniej, a jej pierwiastek \dash{} \textbf{rozrzut} \dash{} ma te same
jednostki co wielkość, przez co właśnie jego da się sobie wyobrazić.
Program~\ref{prog:P25} wraca do szumu z partii jako do przykładu, gdy maszyneria
już istnieje.
\end{rigourbox}
\end{fr}

% ============================================================
\section{Jak bardzo zaszumiony, dokładnie}
\label{sec:P21-variance}
\index{gradient!skala szumu}

\begin{fr}
Populacja ma tu $\val{p21.noise.pop}$ przykładów i rozrzut
$\val{p21.noise.sd}$. Losowane są z niej partie po $B$,
$\val{p21.noise.trials}$ razy przy każdym rozmiarze, i mierzony jest rozrzut
średnich z partii.

Przed liczbami: partia o rozmiarze $B$ uśrednia $B$ niezależnych losowań. Co
powinno stać się z rozrzutem, gdy $B$ rośnie, i jak szybko?
\dotline
\nextframe
\end{fr}

\begin{fr}
\begin{ansblock}
Powinien spadać, a \emph{wariancja} powinna spadać jak $1/B$.
\end{ansblock}

I tak się dzieje. Dla $B$ od $\val{p21.noise.b.lo}$ do $\val{p21.noise.b.hi}$
mierzona wariancja podąża za $\sigma^{2}/B$ z dokładnością do
$\val{p21.noise.tol}$ procent, i skrypt sprawdza to zamiast jakiejkolwiek
pojedynczej liczby \dash{} wariancja zmierzona z czterech tysięcy prób sama
jest wielkością losową i poruszyłaby się razem z ziarnem.

Rozrzut, który naprawdę widać, jest pierwiastkiem z tego, więc spada z
$\val{p21.noise.sd1}$ przy $B = \val{p21.noise.b.lo}$ do
$\val{p21.noise.sd256}$ przy $B = \val{p21.noise.b.hi}$ \dash{} o czynnik
$\val{p21.noise.drop}$ za $\val{p21.noise.b.hi}$ razy więcej pracy.
\end{fr}

\begin{fr}
Ten czynnik jest tym, który ludzie źle odczytują, więc powoli. Przejście z
$B = 1$ do $B = \val{p21.noise.b16}$ to szesnaście razy więcej pracy na krok.

O jaki czynnik spada \emph{rozrzut}?
\nextframe
\end{fr}

\begin{fr}
\ans{$\val{p21.noise.spread16}$}
Cztery, a nie szesnaście, bo rozrzut jest pierwiastkiem z wariancji.

\textbf{Ten pierwiastek jest źródłem \enquote{malejących zysków z rozmiaru
partii}} i jest dokładny, a nie przybliżony. Podwojenie partii kupuje czynnik
$\sqrt{2}$ w szumie za czynnik $2$ w obliczeniach, przy każdym rozmiarze
partii, na zawsze. Nic o modelu nie wchodzi do tego zdania.

\mermaidfig{p21-noise-and-batch}{Dwie liczby opisujące jedną wielkość.}{rozmiar
partii i rozrzut}

Jeszcze jedno przed następną sekcją, bo ona się na tym opiera. To jest rozrzut
\emph{oszacowania}. Ile wynosi rozrzut \emph{kroku}, który wykonujesz?
\nextframe
\end{fr}

\begin{fr}
\begin{ansblock}
$\eta$ razy tyle, bo krok to $\eta$ razy oszacowanie.
\end{ansblock}

Dwa pokrętła, jedna wielkość \dash{} i żadne z nich nie ma znaczenia bez
drugiego, co jest dokładnie tym kształtem, który Program~\ref{prog:P20} znalazł
dla konwencji momentum.
\end{fr}

% ============================================================
\section{Dlaczego zdarza się ogromny krok}
\label{sec:P21-clipping}
\index{gradient!obcinanie}

\begin{fr}
Program~\ref{prog:F06} rozłożył obcinanie gradientu na części: obcinanie po
wartości ucina każdą składową i obraca wektor, obcinanie po normie skaluje go i
nie obraca, a obie nazwy dzieli jedno słowo w każdym API. Powiedział potem, że
pytanie, \emph{dlaczego w ogóle zdarza się ogromny krok}, należy tutaj, bo
potrzebuje modelu szumu.

Model szumu już masz. Powiedz dlaczego.
\dotline
\nextframe
\end{fr}

\begin{fr}
\begin{ansblock}
Bo średnia z partii ma rozrzut, więc partia będzie od czasu do czasu daleko od
gradientu populacji \dash{} a \enquote{od czasu do czasu} jest dokładnie tym,
co rozrzut opisuje.
\end{ansblock}

Nic nie jest nie tak z modelem, danymi ani optymalizatorem, gdy to się zdarza.
To ogon rozkładu, który metoda ma z konstrukcji, a jego częstość wyznaczają
rozrzut i rozmiar partii, i nic więcej.

Czyni to z progu decyzję o rozkładzie, a nie stałą. Zmierzone na tych samych
partiach: ich typowy rozmiar to $\val{p21.clip.typical}$, a próg ustawiony tam
obcina $\val{p21.clip.one}$ procent kroków \dash{} co jest tym, co znaczy
\emph{typowy}, i sprawdzeniem, że arytmetyka się zgadza.
\end{fr}

\begin{fr}
Teraz dwa końce, a przepaść między nimi jest sednem sekcji.

\begin{center}
\begin{tabular}{ll}
\toprule
próg na połowie typowego rozmiaru & obcina $\val{p21.clip.half}$ procent kroków \\
próg na czterokrotności & obcina $\val{p21.clip.four.n}$ z $\val{p21.clip.trials}$ \\
\bottomrule
\end{tabular}
\end{center}

$\val{p21.clip.half}$ procent wobec $\val{p21.clip.four.n}$ z
$\val{p21.clip.trials}$. Co algorytm robi w pierwszym wierszu?
\nextframe
\end{fr}

\begin{fr}
\begin{ansblock}
Podąża za \emph{kierunkiem} gradientu z krokiem ustalonym ręcznie i całkowicie
odrzuca jego długość.
\end{ansblock}

Jest to inny algorytm niż ten, który ktokolwiek sądzi, że uruchamia, a nic tej
różnicy nie zgłasza. Program~\ref{prog:F06} mówi to we własnej ramce; ten
program dodaje, że \textbf{próg ma znaczenie wyłącznie wobec zmierzonego
rozkładu rozmiarów gradientu}, a ten rozkład jest własnością rozmiaru partii
tak samo jak danych.

Pytaniem do progu obcinania nie jest więc \enquote{czy to rozsądna liczba},
tylko \enquote{jaki ułamek kroków obcina}. Pierwsze jest zgadywaniem, drugie
pomiarem, a drugie to dwa wiersze logowania.

Który z dwóch powyższych progów jest bezpieczniejszy w użyciu?
\nextframe
\end{fr}

\begin{fr}
\begin{ansblock}
Ten, który prawie nigdy nie działa.
\end{ansblock}

\begin{trapbox}
\textbf{\enquote{Obcinanie to siatka bezpieczeństwa, więc ciaśniejsza jest
bezpieczniejsza.}} Próg poniżej typowego rozmiaru gradientu nie jest siatką,
tylko przeprojektowaniem: zamienia optymalizator w taki, który używa kierunku
gradientu i ignoruje jego długość, i robi to na niemal każdym kroku, a nie od
czasu do czasu.

Bezpieczne odczytanie jest odwrotnością intuicyjnego. Próg, który działa
rzadko, wykonuje zadanie, do którego został dodany; taki, który działa ciągle,
po cichu zastąpił metodę.
\end{trapbox}
\end{fr}

% ============================================================
\section{Dwie reguły utrzymujące różne rzeczy}
\label{sec:P21-scaling}
\index{optymalizator!reguła liniowego skalowania}

\begin{fr}
Każdy, kto zmieniał rozmiar partii, spotkał to pytanie, a odpowiada się na nie
zwykle regułą, a nie argumentem. Weź najpierw argument.

Aktualizacja to $\eta$ razy oszacowanie. Sekcja 2 podaje wariancję oszacowania
jako $\sigma^{2}/B$. Zapisz wariancję \emph{aktualizacji}.
\dotline
\nextframe
\end{fr}

\begin{fr}
\begin{ansblock}
\[ \Var(\eta\,\hat g) \;=\; \frac{\eta^{2}\sigma^{2}}{B} \]
\end{ansblock}

Jeden wiersz i rozstrzyga pytanie mechanicznie. Pomnóż partię przez $k$ i
zapytaj, co zrobić z $\eta$, żeby ta wielkość się nie ruszyła: $\eta^{2}/B$ ma
zostać bez zmian, więc $\eta$ rośnie o $\sqrt{k}$.

Sprawdzone przy $k = \val{p21.scale.k}$: reguła pierwiastkowa zostawia
wariancję aktualizacji dokładnie na $\val{p21.scale.sqrt}$ raza tym, czym była,
a \textbf{reguła liniowa mnoży ją przez $\val{p21.scale.linear}$}.
\end{fr}

\begin{fr}
Jest to niezręczne, bo regułą faktycznie zalecaną jest liniowa.

Zanim pójdziesz dalej \dash{} nie jest to sprzeczność, a powód mieści się w
jednym słowie. Co utrzymuje stałym reguła liniowa, jeśli nie szum w
aktualizacji?
\nextframe
\end{fr}

\begin{fr}
\begin{ansblock}
Drogę przebytą na \emph{przykład}, a nie na krok. Przy $k$ razy większej partii
robisz $k$ razy mniej kroków, a $k$ razy większy krok pokonuje tyle samo
terenu.
\end{ansblock}

Obie reguły są więc odpowiedziami na dwa różne pytania i obie są dokładnie
poprawne co do własnego. Reguła pierwiastkowa utrzymuje szum każdej
aktualizacji, a liniowa utrzymuje postęp na epokę i pozwala szumowi rosnąć.

\begin{warning}
\textbf{Twierdzenie empiryczne nie jest rozstrzygnięte i ta książka go nie
rozstrzyga.} Który niezmiennik znaczy więcej, zależy od modelu, etapu treningu
i optymalizatora, zgłaszane świadectwa dla reguły liniowej przychodzą z
rozgrzewką i z sufitem, powyżej którego przestaje działać, a spór trwa.

To, co napisano wyżej, jest arytmetyką i obowiązuje niezależnie: to są te dwie
wielkości, a reguła utrzymuje jedną z nich. Ktokolwiek mówi ci, której użyć,
stawia twierdzenie empiryczne, i wolno zapytać, na czym się ono opiera.
\end{warning}

Jest więc trzeci wybór i to on jest robiony bez zastanowienia. Zmniejsz partię
o połowę i zostaw $\eta$: którą z dwóch wielkości to utrzymuje stałą?
\nextframe
\end{fr}

\begin{fr}
\begin{ansblock}
Żadnej.
\end{ansblock}

\begin{trapbox}
\textbf{\enquote{Zmniejszyłem partię o połowę, więc zostawię krok.}} Jest to
trzecia reguła i jedyna z trzech, za którą nikt nie argumentuje: pozwala
szumowi w każdej aktualizacji urosnąć o $\sqrt{2}$ \emph{i} połowi teren
pokonany na przykład.

Użytecznym nawykiem nie jest zapamiętanie reguły. Jest nim zauważenie, że
$\eta$ i $B$ występują razem w jednym wyrażeniu, więc zmiana jednego bez
drugiego jest zmianą wielkości, na którą nie spojrzałeś.
\end{trapbox}
\end{fr}

% ============================================================
\section{Co ukrywa wygładzona krzywa}
\label{sec:P21-smoothing}
\index{średnia bieżąca!a krzywa straty}

\begin{fr}
Surową krzywą straty złożoną z zaszumionych oszacowań trudno się czyta, więc
każdy panel treningowy ją wygładza, a wygładzaczem jest wykładnicza średnia
bieżąca z Programu~\ref{prog:F04}, tyle że ze stratą, a nie z gradientem.

Przy $\beta = \val{p21.ema.beta}$ jej czas połowicznego zaniku wynosi
$\val{p21.ema.halflife}$ kroków, co jest zatwierdzoną liczbą z F04 i jest
wobec niej bramkowane. Załóżmy teraz, że strata naprawdę spada, w jednym kroku,
i zostaje na dole. Co pokazuje wygładzona krzywa?
\dotline
\nextframe
\end{fr}

\begin{fr}
\begin{ansblock}
Połowę spadku po mniej więcej czasie połowicznego zaniku, a resztę przez kilka
następnych.
\end{ansblock}

Zmierzone: wygładzona krzywa przecina punkt środkowy $\val{p21.ema.lag}$ kroków
po zmianie, wobec czasu połowicznego zaniku $\val{p21.ema.halflife}$.
Opóźnienie \textbf{jest} czasem połowicznego zaniku \dash{} nie jest własnością
zmiany i nie zależy od tego, jak duży był spadek.

Wygładzanie nie sprawia więc tylko, że krzywa ładniej wygląda. Przesuwa każde
zdarzenie później o znaną wielkość, a wielkością tą jest liczba, którą wybrałeś,
ustawiając wygładzanie.
\end{fr}

\begin{fr}
\begin{trapbox}
\textbf{\enquote{Strata przestała spadać koło kroku 4000.}} Przy wygładzaniu o
czasie połowicznego zaniku $h$ przestała spadać koło kroku $4000 - h$, a jeśli
porównujesz dwa przebiegi wygładzone inaczej, porównujesz dwa różne opóźnienia.

Działa to i w drugą stronę: interwencja zdaje się działać późno, więc zasługę
przypisuje się temu, co przyszło po niej. Czytaj surową krzywą, gdy pytaniem
jest \emph{kiedy}, a wygładzoną, gdy pytaniem jest \emph{czy}.
\end{trapbox}
\end{fr}

% ============================================================
\section{Jeden mianownik, dwie odpowiedzi}
\label{sec:P21-accumulation}
\index{gradient!akumulacja}

\begin{fr}
Akumulacja gradientu przepuszcza kilka mikropartii, dodaje ich gradienty i
robi jeden krok \dash{} tak trenuje się partię większą niż pamięć. Twierdzenie,
na którym się opiera, mówi, że jest to identyczne z jedną dużą partią.

Program~\ref{prog:F04} już kazał ci spotkać kształt zarzutu. Trzy mikropartie
zawierają $1000$, $10$ i $500$ prawdziwych tokenów i wracają ze średnimi
stratami $2$, $8$ i $3$. Ile wynosi średnia strata na
$\val{p21.acc.tokens}$ tokenach?
\dotline
\nextframe
\end{fr}

\begin{fr}
\ans{$\val{p21.acc.pooled}$}
A tym, co liczy akumulacja, jeśli każda mikropartia już podzieliła przez własną
liczbę tokenów, jest średnia z $2$, $8$ i $3$ \dash{}
$\val{p21.acc.accumulated}$, czyli $\val{p21.acc.ratio}$ raza więcej.

\textbf{Różnica nie jest w raporcie. Jest w gradiencie}, bo strata, którą
uśredniono źle, jest stratą, którą zróżniczkowano. Przebieg optymalizuje cel, w
którym mikropartia z dziesięcioma tokenami liczy się tyle samo, co ta z
tysiącem.
\end{fr}

\begin{fr}
Rodzi to pytanie, przez które warto poświęcić na to sekcję, a nie ostrzeżenie:
dlaczego ktokolwiek to wypuszcza?

Policz te same dwie liczby dla trzech mikropartii zawierających po $500$
tokenów.
\nextframe
\end{fr}

\begin{fr}
\begin{ansblock}
Są równe \dash{} dokładnie, na ułamkach.
\end{ansblock}

\textbf{Dlatego to się ukrywa.} Przy równych liczbach tokenów oba wyrażenia są
tym samym wyrażeniem, więc test napisany na dopełnionych danych przechodzi,
mały przebieg z jednorodnymi partiami przechodzi, a wada pojawia się dopiero
przy nierównych mikropartiach \dash{} czyli na prawdziwym tekście, przy skali,
w przebiegu, który się liczy.

Jest to \enquote{bezpieczne dla danych, których nie próbowałeś} z
Programu~\ref{prog:P02} po raz trzeci w tej książce, i tożsamość rozgłaszania z
Programu~\ref{prog:P07} po raz drugi: wada niewidoczna na danych testowych i
rozstrzygająca na korpusie.
\end{fr}

\begin{fr}
\begin{trapbox}
\textbf{\enquote{Akumulacja po czterech mikropartiach to partia czterokrotnie
większa.}} Tylko wtedy, gdy cztery zawierają tyle samo prawdziwych tokenów. W
przeciwnym razie jest to średnia ze średnich tam, gdzie chciano średniej z
puli, i dzieje się po cichu: żadnego błędu, żadnego ostrzeżenia i krzywa straty
wyglądająca zupełnie normalnie, bo jest krzywą innego celu.

Naprawą jest ta, którą podaje Program~\ref{prog:F04} \dash{} nieś licznik i
mianownik osobno i podziel raz, na końcu \dash{} a testem ten, który ta sekcja
podpowiada: uruchom przypadek z celowo nierównymi mikropartiami i porównaj
gradient akumulowany z gradientem dużej partii.
\end{trapbox}
\end{fr}

% ============================================================
\section{Dwa sposoby przez rzut monetą}
\label{sec:P21-estimators}
\index{gradient!estymator funkcji wynikowej}
\index{gradient!reparametryzacja}

\begin{fr}
Wszystko do tej pory dotyczyło szumu, którego nie chciałeś. Ta sekcja dotyczy
losowości będącej częścią modelu: kroku losowania \emph{wewnątrz} tego, co jest
różniczkowane.

Nie jest to nic egzotycznego. Polityka losuje akcję, model wariacyjny losuje
zmienną ukrytą, dekoder losuje token. W każdym przypadku wielkością do
maksymalizacji jest średnia po tym, co może zostać wylosowane, i chce się jej
gradientu.

Powiedz, co czyni to trudnym.
\dotline
\nextframe
\end{fr}

\begin{fr}
\begin{ansblock}
Krok losowania nie ma pochodnej. Nie da się zróżniczkować \enquote{losowania}.
\end{ansblock}

Są na to dwa sposoby i różnią się całkowicie charakterem. Ta sekcja nie
wyprowadza żadnego z nich \dash{} oba to wiersz algebry, który Programy
\ref{prog:P24} i \ref{prog:P26} podadzą lepiej \dash{} i robi zamiast tego to,
co ta książka potrafi: mierzy, co je dzieli.

\mermaidfig{p21-two-estimators}{Dwie drogi i koszt wyboru.}{dwie drogi przez
losowanie}
\end{fr}

\begin{fr}
\textbf{Estymator funkcji wynikowej}, czyli to, czym jest gradient polityki:
pomnóż uzyskaną wartość przez pochodną logarytmu prawdopodobieństwa jej
uzyskania. Potrzebuje \emph{wartości} funkcji i niczego więcej \dash{} funkcja
może być czarną skrzynką, symulatorem, oceną człowieka.

\textbf{Sztuczka z reparametryzacją}, czyli to, czego używa autoenkoder
wariacyjny: zapisz losowanie jako ustaloną liczbę losową przepuszczoną przez
różniczkowalną funkcję parametrów i zróżniczkuj wprost przez nią. Potrzebuje
różniczkowalnej drogi od parametru do wartości.

Oba są nieobciążone. Powiedz, co zostaje do wyboru między nimi.
\nextframe
\end{fr}

\begin{fr}
\begin{ansblock}
Czy droga w ogóle istnieje \dash{} a gdy istnieje, \textbf{wariancja}.
\end{ansblock}

Zmierzone na jednym zadaniu, gdzie oba działają, przy
$\val{p21.grad.samples}$ próbkach i prawdziwym gradiencie
$\val{p21.grad.true}$:

\begin{center}
\begin{tabular}{lll}
\toprule
wymiar & funkcja wynikowa & reparametryzacja \\
\midrule
$\val{p21.grad.d.lo}$ & $\val{p21.grad.score.lo}$ & $\val{p21.grad.repar}$ \\
$\val{p21.grad.d.hi}$ & $\val{p21.grad.score.hi}$ & $\val{p21.grad.repar}$ \\
\bottomrule
\end{tabular}
\end{center}

Przeczytaj najpierw prawą kolumnę. Co ona zrobiła?
\nextframe
\end{fr}

\begin{fr}
\begin{ansblock}
Nic. Wariancja estymatora reparametryzowanego jest taka sama przy
$\val{p21.grad.d.hi}$ wymiarach jak przy $\val{p21.grad.d.lo}$.
\end{ansblock}

Dokładnie taka sama i nie przez przypadek: estymator jednej składowej jest
funkcją tej składowej i niczego więcej, więc pozostałe $99$ wymiarów do niego
nie wchodzi. Skrypt sprawdza to przy każdym wymiarze.

Wariancja estymatora funkcji wynikowej idzie tymczasem z
$\val{p21.grad.score.lo}$ do $\val{p21.grad.score.hi}$, czyli o czynnik
$\val{p21.grad.ratio.hi}$ przy $\val{p21.grad.d.hi}$ wymiarach \dash{} bo mnoży
przez wartość \emph{całej} funkcji, a cała funkcja rośnie z wymiarem.
\end{fr}

\begin{fr}
Teraz konsekwencja, czyli zdanie, dla którego istnieje ta sekcja. Oba
estymatory są nieobciążone. Po $\val{p21.grad.samples}$ próbkach przy
$\val{p21.grad.d.hi}$ wymiarach oszacowanie reparametryzowane jest o
$\val{p21.grad.off.repar}$ procent od prawdziwej wartości.

O ile odbiega oszacowanie funkcji wynikowej?
\nextframe
\end{fr}

\begin{fr}
\ans{$\val{p21.grad.off.score}$ procent}
$\val{p21.grad.score.est}$ wobec prawdziwego $\val{p21.grad.true}$, po
czterdziestu tysiącach próbek, z estymatora, który jest średnio dokładnie
poprawny.

\textbf{Nieobciążony to nie to samo, co użyteczny.} To cała rozwidlenie: model
wariacyjny ma różniczkowalną drogę i z niej korzysta, a metoda gradientu
polityki jej nie ma i za to płaci \dash{} dlatego literatura o gradientach
polityki jest w dużej mierze literaturą o redukcji wariancji, a literatura o
autoenkoderach wariacyjnych nie.
\end{fr}

\begin{fr}
\begin{trapbox}
\textbf{\enquote{Oba estymatory są nieobciążone, więc są wymienne.}} Szacują tę
samą wielkość i nie są wymienne, bo estymatora używa się skończenie wiele razy.
Istotnym porównaniem nie jest to, czy średnia jest poprawna, lecz jak daleko od
niej zostawia cię realistyczna liczba próbek, a to wyznacza wariancja.

Jest to to samo rozróżnienie, co w sekcji 1, na drugim końcu programu: tam
nieobciążoność nie mówiła nic o jednej partii, tutaj nie mówi nic o jednym
przebiegu.
\end{trapbox}
\end{fr}

\begin{fr}
To zamyka program i warto powiedzieć, co obowiązywało przez cały czas.

Każde zażalenie, którym ten program się zajął \dash{} zaszumiona krzywa,
ogromny krok, krok, który się nie przeniósł, wygładzona krzywa kłamiąca o tym,
kiedy, akumulowana partia, która nią nie była, estymator, który nie chciał
zbiec \dash{} było stwierdzeniem o \emph{wariancji} i żadne nie było
stwierdzeniem o modelu.

Warto to nieść, bo odruchem w każdym z tych przypadków jest spojrzeć na model.
Nawykiem, dla którego istnieje ten program, jest zapytać najpierw, jaka wielkość
jest szacowana, z ilu próbek i z jakim rozrzutem \dash{} a odpowiedzią jest
zwykle dwuwierszowa arytmetyka, której nikt nie zapisał.
\end{fr}

\begin{summarybox}
\sumitem{1--2}{\result{Gradient z partii jest \textbf{oszacowaniem} gradientu, którego
  chciałeś. Szum jest metodą, a nie wadą w niej}.}
\sumitem{3--4}{\result{\emph{Nieobciążoność} znaczy, że średnia po każdej partii, jaka
  mogła zostać wylosowana, jest dokładna \dash{} udowodnione tu na wszystkich
  $\val{p21.pop.subsets}$ z nich, na ułamkach, co jest dowodem dla tej
  populacji, a nie świadectwem o niej}.}
\sumitem{5--6}{\result{Nie mówi nic o jednym losowaniu: średnie z partii biegną od
  $\val{p21.pop.lo}$ do $\val{p21.pop.hi}$ wokół $\val{p21.pop.mean}$, a jedna
  z nich ma zły znak}.}
\sumitem{7}{Zależność w przód jest zadeklarowana: sama wariancja należy do
  Programu~\ref{prog:P24}, a koncentracja do Programu~\ref{prog:P25}, dwie
  części później, i P25 wraca do szumu z partii, gdy maszyneria już istnieje.}
\sumitem{8--9}{\result{Wariancja średniej z partii spada jak $1/B$, zmierzone z
  dokładnością do $\val{p21.noise.tol}$ procent dla $B$ od
  $\val{p21.noise.b.lo}$ do $\val{p21.noise.b.hi}$}.}
\sumitem{10--11}{\result{\textbf{Rozrzut spada jak $1/\sqrt{B}$}, więc
  $\val{p21.noise.b16}$ razy większa partia kupuje $\val{p21.noise.spread16}$
  razy mniej szumu. Ten pierwiastek jest całością \enquote{malejących zysków} i
  jest dokładny, a nie przybliżony}.}
\sumitem{12}{\result{Rozrzut kroku to $\eta$ razy rozrzut oszacowania, więc $\eta$ i
  $B$ to dwie liczby opisujące jedną wielkość}.}
\sumitem{13--14}{\result{Ogromny krok zdarza się, bo średnia z partii ma rozrzut. Nic
  nie jest nie tak z modelem, danymi ani optymalizatorem, gdy się zdarza}.}
\sumitem{15--17}{\result{Próg obcinania to decyzja o rozkładzie: na połowie typowego
  rozmiaru obcina $\val{p21.clip.half}$ procent kroków, a na czterokrotności
  $\val{p21.clip.four.n}$ z $\val{p21.clip.trials}$. \textbf{Pierwszy jest
  innym algorytmem}, podążającym za kierunkiem i odrzucającym długość, a
  bezpieczny próg to ten, który prawie nigdy nie działa}.}
\sumitem{18--19}{\result{Wariancja aktualizacji to $\eta^{2}\sigma^{2}/B$, co
  rozstrzyga pytanie o skalowanie mechanicznie}.}
\sumitem{20--21}{\result{Skalowanie $\eta$ przez $\sqrt{k}$ utrzymuje tę wariancję
  \textbf{dokładnie} stałą; skalowanie przez $k$ mnoży ją przez
  $\val{p21.scale.linear}$ i utrzymuje zamiast tego teren pokonany na przykład.
  Dwie reguły, dwa niezmienniki, a pytanie empiryczne nie jest tu
  rozstrzygnięte}.}
\sumitem{22}{\result{Zostawienie $\eta$ przy zmianie $B$ nie utrzymuje żadnej z nich
  stałej, i dlatego jest jedynym wyborem, za którym nikt nie argumentuje}.}
\sumitem{23--24}{\result{Wygładzona krzywa straty opóźnia się o swój czas połowicznego
  zaniku: $\val{p21.ema.lag}$ kroków zmierzone wobec $\val{p21.ema.halflife}$
  zatwierdzonego przez Program~\ref{prog:F04}. Opóźnienie nie zależy od
  wielkości zmiany}.}
\sumitem{25}{\result{Czytaj surową krzywą, gdy pytaniem jest \emph{kiedy}, a wygładzoną,
  gdy jest nim \emph{czy}}.}
\sumitem{26--27}{\result{Akumulacja mikropartii o różnych liczbach tokenów daje
  $\val{p21.acc.accumulated}$ tam, gdzie duża partia daje
  $\val{p21.acc.pooled}$ \dash{} i jest to \textbf{gradient}, a nie raport, bo
  strata uśredniona źle jest stratą zróżniczkowaną}.}
\sumitem{28--30}{\result{Przy równych liczbach tokenów oba są identyczne na ułamkach, i
  dokładnie dlatego to wychodzi na produkcję: jest niewidoczne na dopełnionych
  danych i rozstrzygające na prawdziwym tekście}.}
\sumitem{31--33}{\result{Krok losowania nie ma pochodnej, a są dwa sposoby, by go
  obejść: funkcja wynikowa, która potrzebuje samej wartości, i
  reparametryzacja, która potrzebuje różniczkowalnej drogi}.}
\sumitem{34--35}{\result{Oba są nieobciążone. Wariancja estymatora reparametryzowanego
  to $\val{p21.grad.repar}$ przy \emph{każdym} wymiarze; funkcji wynikowej
  dochodzi do $\val{p21.grad.score.hi}$, o czynnik $\val{p21.grad.ratio.hi}$}.}
\sumitem{36--38}{\result{Po $\val{p21.grad.samples}$ próbkach oszacowanie funkcji
  wynikowej wciąż odbiega o $\val{p21.grad.off.score}$ procent.
  \textbf{Nieobciążony to nie to samo, co użyteczny}, i to jest rozwidlenie
  między modelem wariacyjnym a gradientem polityki}.}
\sumitem{39}{Każde zażalenie w tym programie było stwierdzeniem o
  \textbf{wariancji} i żadne nie było stwierdzeniem o modelu.}
\end{summarybox}

\canyou

\begin{testexercises}
\item Partia o rozmiarze $32$ daje oszacowania gradientu o rozrzucie
  $\num{0.40}$. Jaki rozmiar partii połowiłby ten rozrzut i ile to kosztuje?
  \answerto{$128$. Rozrzut spada jak $1/\sqrt{B}$, więc jego połowienie
  wymaga czterokrotnie większej partii \dash{} i czterokrotnie więcej obliczeń
  na krok.}
\item Przebieg przy partii $256$ używa $\eta = \num{3e-4}$. Przechodzisz na
  partię $1024$. Ile wynosi $\eta$ według reguły liniowej, a ile według
  pierwiastkowej?
  \answerto{$k = 4$, więc liniowa daje $\num{1.2e-3}$, a pierwiastkowa
  $\num{6e-4}$. Pierwsza utrzymuje stały teren pokonany na przykład i mnoży
  szum każdej aktualizacji przez $4$; druga utrzymuje ten szum dokładnie
  stały.}
\item Panel wygładza stratę przy $\beta = \num{0.98}$. Zmiana została
  wprowadzona w kroku $1000$. Mniej więcej kiedy wygładzona krzywa pokaże
  połowę jej efektu?
  \answerto{Około kroku $1034$. Czas połowicznego zaniku to
  $\ln(\num{0.5})/\ln(\num{0.98}) \approx 34$ kroków, a opóźnienie jest tym
  czasem niezależnie od wielkości zmiany.}
\item Dwie mikropartie zawierają $600$ i $200$ prawdziwych tokenów ze średnimi
  stratami $\num{1.5}$ i $\num{3.5}$. Co zgłasza akumulacja, a co powinna?
  \answerto{Akumulacja zgłasza średnią ze średnich, $\num{2.5}$. Średnia z puli
  to $(600 \times \num{1.5} + 200 \times \num{3.5})/800 = \num{2.0}$. Gradient
  jest błędny o ten sam stosunek, nie tylko liczba na ekranie.}
\item Trenujesz model, którego wyjście trafia do nieróżniczkowalnego oceniacza.
  Który z dwóch estymatorów jest dostępny i czego się spodziewać?
  \answerto{Tylko estymator funkcji wynikowej, bo nie ma różniczkowalnej drogi
  przez oceniacz. Spodziewaj się gradientu o dużej wariancji i spodziewaj się
  wysiłku włożonego w redukcję wariancji \dash{} do tego właśnie służą wartości
  bazowe i oszacowania przewagi.}
\end{testexercises}

\begin{furtherproblems}
\item Pokaż, że jeśli $\hat g$ jest średnią z $B$ niezależnych losowań o
  wariancji $\sigma^{2}$, to $\Var(\hat g) = \sigma^{2}/B$, i powiedz, które
  słowo w tym zdaniu wykonuje pracę.
  \answerto{Wariancja sumy niezależnych wielkości jest sumą ich wariancji, więc
  suma $B$ losowań ma wariancję $B\sigma^{2}$; dzielenie przez $B$ dzieli
  wariancję przez $B^{2}$, zostawiając $\sigma^{2}/B$. Pracę wykonuje słowo
  \emph{niezależnych}: losowania skorelowane tego nie spełniają, i dlatego
  losowanie partii z posortowanego albo poklastrowanego zbioru łamie model
  szumu, a nie jest tylko niechlujne.}
\item Zapisz, co utrzymuje stałym każda z dwóch reguł skalowania, jako
  wyrażenie, i powiedz, co utrzymuje stałym trzeci częsty wybór \dash{}
  zostawienie $\eta$.
  \answerto{Skalowanie pierwiastkowe utrzymuje stałe $\eta^{2}/B$, czyli
  wariancję jednej aktualizacji. Skalowanie liniowe utrzymuje stały teren
  pokonany na przykład, bo $k$ razy większa partia oznacza $k$ razy mniej
  kroków. Zostawienie $\eta$ nie utrzymuje stałego żadnego z nich: wariancja
  aktualizacji zmienia się o $1/k$, a teren na przykład również o $1/k$, i
  dlatego jest to wybór, za którym nikt nie argumentuje, a nie trzecia reguła.}
\item Kolega zgłasza, że obcinanie \enquote{nic nie robi}, bo jego usunięcie
  prawie nie zmienia krzywej straty. Jaki pomiar rozstrzygnąłby, czy to prawda,
  i co znaczyłaby każda odpowiedź?
  \answerto{Zaloguj ułamek kroków, które są obcinane. Jeśli jest bliski zeru,
  obcinanie faktycznie nic nie robi w tym przebiegu, a próg spełnia swoją rolę
  zabezpieczenia. Jeśli jest duży, obcinanie działa na większości kroków,
  optymalizator podąża za kierunkiem gradientu z ręcznie ustaloną długością, a
  \enquote{usunięcie prawie nic nie zmienia} znaczy, że w długościach było mało
  informacji \dash{} co jest odkryciem o wiele ciekawszym i zupełnie innym.}
\item Wariancja estymatora reparametryzowanego nie zmieniła się między jednym a
  stu wymiarami. Wyjaśnij dlaczego i powiedz, co by to zepsuło.
  \answerto{Estymator jednej składowej parametru jest funkcją losowania tej
  składowej i niczego więcej, więc pozostałe wymiary nigdy do niego nie
  wchodzą. Zepsułoby to sprzężenie składowych przez uśrednianą funkcję w taki
  sposób, że pochodna po jednej zależałaby od wszystkich \dash{} co jest
  częste, i dlatego ta płaskość jest własnością tego przykładu, a nie ogólnym
  prawem. Ogólny jest kontrast: estymator funkcji wynikowej mnoży przez wartość
  całej funkcji, więc podchwytuje wymiar niezależnie od tego, jaka to funkcja.}
\item Oba estymatory są nieobciążone, a jeden z nich jest bezużyteczny przy stu
  wymiarach. Pogódź te dwa zdania w jednym zdaniu i nazwij wielkość, która je
  godzi.
  \answerto{Nieobciążoność jest stwierdzeniem o granicy nieskończenie wielu
  próbek, a estymatora używa się skończenie wiele razy, więc o użyteczności
  rozstrzyga \emph{wariancja} \dash{} to ona wyznacza, jak daleko realistyczna
  próba zostawia cię od wartości, którą dałaby granica.}
\end{furtherproblems}
